\documentclass[11pt]{article}
\usepackage[margin=0.8in]{geometry}
\usepackage[]{xcolor}
\usepackage{amsmath}
\parindent=0pt

\title{Memory Management cont.}
\author{gordonbchen}
\date{}

\begin{document}
\maketitle


\section{Allocation policy}

Local allocation policy: only page in same proc

Global allocation: takes more time to scan, but less page faults
\begin{itemize}
    \item better when process size chagnes (big $\to$ small)
\end{itemize}

Allocation 
\begin{itemize}
    \item allocate equally
    \item proportional to process size
    \item allocate based on page fault frequency: track page fault freq.

        upper and lower bound for page faults, allocate and deallocate pages

    \item working set page replacement alg only works for local allocation b/c of virtual times

    \item swapping processes to disk
\end{itemize}

page size
\begin{itemize}
    \item on avg, half of last page is empty
    \item smaller page = less wasted space, but larger table, slower loading time
    \item s = avg proc size, p = page size, e = page table entry size

        overhead = (s / p) e + p/2

        optimal page size = sqrt(2 s e)
\end{itemize}

separate instruction and data spaces
\begin{itemize}
    \item I: instruction
    \item D: data
    \item fork doesn't have to re-copy instruction, only copy data space on write

        process table points to same program/data space

        data space is marked R, after copy on write, mark R/W
    \item DLL: shared libraries

        compiling printf: don't have duplicate copies of the library. implementation of shared library
            is just pointer, not actual binary in compiled program.
\end{itemize}

mapped files: map RAM address to IO.

cleaning
\begin{itemize}
    \item paging daemon: in backgroup, periodically checks available memory

        if available mem too low, pre-select pages to evice, if dirty then schedule writes

    \item 2-handed clock: first hand is cleaning daemon, 2nd hand is page replacement

        need enough buffer b/t cleaning and replacment
\end{itemize}

instruction can span multiple bytes,autoincreementing must have handware support

A starts IO request, gets suspeneded. B runs and get a page fault. Must not be page A is using for IO.
\begin{itemize}
    \item pinning "locking" a page into memory
    \item kenerl buffers for IO, then copy to use page frame
\end{itemize}

backing store: 1 to 1 map or disk map that maps from virtual page to swap area on HDD

segmentation
\begin{itemize}
    \item 1d address space with growing tables can collide
    \item virtual addr space for each table
    \item segments: symbol table, source text, constants, parse tree, call stack...
    \item page table for each process and each segment
\end{itemize}

segmentation with paging: MULTICS
\begin{enumerate}
    \item segment desscription points to page table RAM addr for segment
    \item segment number (18-bit) + page number (6-bit) + page offset (10-bit)

        $2^{18}$ segments, $2^6$ pages, $2^{10}$ byte pages

    \item now TLB needs segment number too (and virtual to page, protection, age, present/dirty)
    \item 256K segments, 64k 36-bit words
\end{enumerate}

Pentium
\begin{enumerate}
    \item segmentation obsolete in x86\_64
    \item 16k segements, 1 billion 32-bit words
    \item each program has LDT (local descriptor table), 1 GDT (global descriptor table) for OS
    \item selector: index + gdt/ldt

        use selector to find descriptor, descriptor base addr + offset

    \item if paging, then go to page table b/c this is a virtual addr

        then 2-level page table: dir $\to$ page $\to$ byte (using offset)
\end{enumerate}

\end{document}
